% GENERATED FILE. Edit canonical lessons, then run npm run generate:books.
% canonical-source-hash: fnv1a64:a78ed6bed53b386c
% canonical-lessons: SA-C193-krayah, SA-C193-vikrayah, SA-C193-vastu, SA-C193-padarthah, SA-C193-upakaranam

\chapter{Purchase, Sale, Object, Item, Tool}
\label{ch:sa-purchase-sale-object-item-tool}
\hlchaptermodality{193}

\begin{chapteropening}
\textbf{By the end of this chapter:} \emph{I can say a purchase, a sale, an object, an item and a tool.}

Complete the last lesson of chapter 193: I can say a purchase, a sale, an object, an item and a tool.
\end{chapteropening}

\section[\texorpdfstring{\sk{क्रयः} (krayaḥ)}{क्रयः (krayaḥ)}]{\sk{क्रयः} (krayaḥ) — a purchase}

\label{lesson:SA-C193-krayah}



\begin{quote}
Before the new one: say the Sanskrit for clothing, then the Sanskrit for a garment.
\end{quote}

\subsection*{You'll want to know: \sk{क्रयः}}
\textbf{\sk{क्रयः}} — \emph{krayaḥ} — \textquotedblleft{}a purchase\textquotedblright{}.

\begin{grammarlens}[title={What you've built}]
One more word for this chapter.
\end{grammarlens}

\subsection*{Your turn}
\begin{itemize}
\raggedright
  \item \emph{Say it:} \emph{krayaḥ}
  \item \emph{Say it:} \emph{krayaḥ}, once more
  \item \emph{Say it:} say \emph{krayaḥ}
  \item \emph{Recall:} say the Sanskrit for clothing, then the Sanskrit for a garment, then say \emph{krayaḥ} again
\end{itemize}

\begin{tcolorbox}[breakable,colback=teal!4,colframe=teal!35!black,title={Before you move on}]
What does \sk{क्रयः} mean? (\textquotedblleft{}A purchase\textquotedblright{}.) Say it once more.
\end{tcolorbox}
\section[\texorpdfstring{\sk{विक्रयः} (vikrayaḥ)}{विक्रयः (vikrayaḥ)}]{\sk{विक्रयः} (vikrayaḥ) — a sale}

\label{lesson:SA-C193-vikrayah}



\begin{quote}
Before the new one: say the Sanskrit for a garment, then the Sanskrit for a purchase.
\end{quote}

\subsection*{You'll want to know: \sk{विक्रयः}}
\textbf{\sk{विक्रयः}} — \emph{vikrayaḥ} — \textquotedblleft{}a sale\textquotedblright{}.

\begin{grammarlens}[title={What you've built}]
One more word for this chapter.
\end{grammarlens}

\subsection*{Your turn}
\begin{itemize}
\raggedright
  \item \emph{Say it:} \emph{vikrayaḥ}
  \item \emph{Say it:} \emph{vikrayaḥ}, once more
  \item \emph{Say it:} say \emph{vikrayaḥ}
  \item \emph{Recall:} say the Sanskrit for a garment, then the Sanskrit for a purchase, then say \emph{vikrayaḥ} again
\end{itemize}

\begin{tcolorbox}[breakable,colback=teal!4,colframe=teal!35!black,title={Before you move on}]
What does \sk{विक्रयः} mean? (\textquotedblleft{}A sale\textquotedblright{}.) Say it once more.
\end{tcolorbox}
\section[\texorpdfstring{\sk{वस्तु} (vastu)}{वस्तु (vastu)}]{\sk{वस्तु} (vastu) — an object, a thing}

\label{lesson:SA-C193-vastu}



\begin{quote}
Before the new one: say the Sanskrit for a purchase, then the Sanskrit for a sale.
\end{quote}

\subsection*{You'll want to know: \sk{वस्तु}}
\textbf{\sk{वस्तु}} — \emph{vastu} — \textquotedblleft{}an object, a thing\textquotedblright{}.

\begin{grammarlens}[title={What you've built}]
One more word for this chapter.
\end{grammarlens}

\subsection*{Your turn}
\begin{itemize}
\raggedright
  \item \emph{Say it:} \emph{vastu}
  \item \emph{Say it:} \emph{vastu}, once more
  \item \emph{Say it:} say \emph{vastu}
  \item \emph{Recall:} say the Sanskrit for a purchase, then the Sanskrit for a sale, then say \emph{vastu} again
\end{itemize}

\begin{tcolorbox}[breakable,colback=teal!4,colframe=teal!35!black,title={Before you move on}]
What does \sk{वस्तु} mean? (\textquotedblleft{}An object, a thing\textquotedblright{}.) Say it once more.
\end{tcolorbox}
\section[\texorpdfstring{\sk{पदार्थः} (padārthaḥ)}{पदार्थः (padārthaḥ)}]{\sk{पदार्थः} (padārthaḥ) — an item, a thing}

\label{lesson:SA-C193-padarthah}



\begin{quote}
Before the new one: say the Sanskrit for a sale, then the Sanskrit for an object.
\end{quote}

\subsection*{You'll want to know: \sk{पदार्थः}}
\textbf{\sk{पदार्थः}} — \emph{padārthaḥ} — \textquotedblleft{}an item, a thing\textquotedblright{}.

\begin{grammarlens}[title={What you've built}]
One more word for this chapter.
\end{grammarlens}

\subsection*{Your turn}
\begin{itemize}
\raggedright
  \item \emph{Say it:} \emph{padārthaḥ}
  \item \emph{Say it:} \emph{padārthaḥ}, once more
  \item \emph{Say it:} say \emph{padārthaḥ}
  \item \emph{Recall:} say the Sanskrit for a sale, then the Sanskrit for an object, then say \emph{padārthaḥ} again
\end{itemize}

\begin{tcolorbox}[breakable,colback=teal!4,colframe=teal!35!black,title={Before you move on}]
What does \sk{पदार्थः} mean? (\textquotedblleft{}An item, a thing\textquotedblright{}.) Say it once more.
\end{tcolorbox}
\section[\texorpdfstring{\sk{उपकरणम्} (upakaraṇam)}{उपकरणम् (upakaraṇam)}]{\sk{उपकरणम्} (upakaraṇam) — a tool, an instrument}

\label{lesson:SA-C193-upakaranam}



\begin{quote}
Before the new one: say the Sanskrit for an object, then the Sanskrit for an item.
\end{quote}

\subsection*{You'll want to know: \sk{उपकरणम्}}
\textbf{\sk{उपकरणम्}} — \emph{upakaraṇam} — \textquotedblleft{}a tool, an instrument\textquotedblright{}.

\begin{grammarlens}[title={What you've built}]
One more word for this chapter.
\end{grammarlens}

\subsection*{Your turn}
\begin{itemize}
\raggedright
  \item \emph{Say it:} \emph{upakaraṇam}
  \item \emph{Say it:} \emph{upakaraṇam}, once more
  \item \emph{Say it:} say \emph{upakaraṇam}
  \item \emph{Recall:} say the Sanskrit for an object, then the Sanskrit for an item, then say \emph{upakaraṇam} again
\end{itemize}

\begin{tcolorbox}[breakable,colback=teal!4,colframe=teal!35!black,title={Before you move on}]
What does \sk{उपकरणम्} mean? (\textquotedblleft{}A tool, an instrument\textquotedblright{}.) Say it once more.
\end{tcolorbox}
